% sec-11: references, abbreviations, availability, back matter.
\section*{Back Matter}
\phantomsection\addcontentsline{toc}{section}{Back Matter}

\bmhead{Abbreviations}

\draftinstr{full-apply builds a full-width two-column abbreviation table.
Include at least: AMP, ARPA-H, CTEP, ctDNA, DLT, DOI, FNIH, FRO, HHMI, IND, IRB,
ISGPS, KRAS, LLM, MTD, NIH, NSF, ODE, OS, PDAC, PFS, PK, PD, QSP, RP2D, SAE,
SBIR, TIP, VVUQ.}

\bmhead{Data and Code Availability}

\draftinstr{One paragraph: every work cited to a DOI is openly deposited; the
repository is at the URL on the cover; simulation code, verification notebooks,
and every figure's TikZ source are in it. State the repository version, v4.4.0.}

\bmhead{Author Contributions and Conflicts}

\draftinstr{One paragraph. Kevin Kawchak is the sole author. He is chief
executive of ChemicalQDevice and holds no equity in Revolution Medicines or any
robotic surgery vendor. No sponsor, contract research organization, or medical
society reviewed this paper. The work was adapted using Claude Code Opus 5.}

\bmhead{Citation}

\draftinstr{full-apply writes the citation block: Kawchak K. 10 Funding
Applications: A Phase 1, First-in-Human, PDAC Clinical Trial Protocol of a
LLM-Directed Robotic Whipple with Daraxonrasib (RMC-6236). Zenodo. 2026;
10.5281/zenodo.xxxxxxxx, hyperlinked to
https://doi.org/10.5281/zenodo.xxxxxxxx. The placeholder is deliberate and is
kept through final-apply.}

\vspace{0.2em}
\apprule{0.5pt}
\vspace{-0.1em}

\bmhead{References}
\vspace{-0.35em}
\begin{apprefs}
\bibliographystyle{unsrt}
\bibliography{references}
\end{apprefs}
